% =============================================================================
% Introduzione: capitolo non numerato in apertura del corpo della tesi, con voce
% nell'indice. Testo da guideline/Intro&Abstract.docx, riportato parola per parola;
% le citazioni, che nel Word erano scritte come numeri della bibliografia, sono \cite.
% =============================================================================
\chapter*{Introduzione}
\label{cap:introduzione}
\addcontentsline{toc}{chapter}{Introduzione}
% Stile "plain": con "fancy" un capitolo senza numero stamperebbe "Capitolo 0".
\pagestyle{plain}

I processi di regolazione rivestono un ruolo centrale nello sviluppo durante la prima infanzia, sostenendo la progressiva capacità del bambino di modulare emozioni, attenzione e comportamento in funzione delle richieste dell'ambiente. Tra i 18 e i 36 mesi tali competenze risultano ancora in fase di organizzazione e consolidamento e si sviluppano attraverso la progressiva integrazione di componenti fisiologiche, emotive, attentive e comportamentali, in stretta relazione con le caratteristiche individuali e con i processi di co-regolazione sostenuti dall'adulto~\cite{calkins2007,blair2015}. La natura evolutiva e multidimensionale del funzionamento regolativo rende questa fase particolarmente rilevante per comprenderne sia le traiettorie tipiche sia le possibili condizioni di vulnerabilità.

Le difficoltà nei processi di regolazione assumono particolare interesse nella psicopatologia dello sviluppo, poiché possono manifestarsi trasversalmente a differenti profili e traiettorie evolutive. In una prospettiva dimensionale e transdiagnostica, la disregolazione non viene quindi considerata esclusivamente in relazione a specifiche categorie diagnostiche, ma come una dimensione del funzionamento che può assumere rilevanza rispetto a molteplici esiti evolutivi~\cite{beauchaine2019,astle2022}. Tale prospettiva risulta particolarmente significativa nei primi anni di vita, quando le competenze regolative sono ancora emergenti e la distinzione tra variabilità evolutiva e possibili segnali di difficoltà richiede una lettura attenta delle caratteristiche e del contesto in cui i comportamenti si manifestano.

La possibilità di individuare e descrivere precocemente tali differenze pone tuttavia rilevanti problemi sul piano della valutazione. L'autoregolazione non è direttamente osservabile, ma viene inferita attraverso comportamenti e prestazioni rilevati in situazioni specifiche; inoltre, nei bambini molto piccoli le diverse componenti del funzionamento regolativo risultano ancora scarsamente differenziate e particolarmente sensibili alle condizioni contestuali. Ne deriva la necessità di adottare procedure adeguate alle caratteristiche evolutive della fascia 18--36 mesi e di integrare differenti modalità e fonti di osservazione, al fine di ottenere una rappresentazione più articolata del funzionamento regolativo e ridurre il rischio di inferenze basate su singole prestazioni~\cite{mcclelland2012,snyder2021}.

In questo quadro si colloca il Baby-FE, una batteria volta ad ampliare le possibilità di osservazione delle competenze regolative e del funzionamento esecutivo emergente attraverso attività strutturate e adeguate alle caratteristiche evolutive dei bambini tra i 18 e i 36 mesi. Il presente lavoro si inserisce all'interno di uno studio più ampio sull'utilizzo dello strumento e si concentra sull'approfondimento di bambini caratterizzati dalla presenza di condizioni di rischio evolutivo, con l'obiettivo di esplorare quanto il Baby-FE possa risultare informativo nella comprensione delle loro caratteristiche regolative. In particolare, l'attenzione è rivolta alla capacità dello strumento di distinguere differenti livelli di prestazione e di cogliere competenze emergenti, considerando congiuntamente il punteggio complessivo, le modalità di risposta alle singole prove e le informazioni provenienti da ulteriori fonti osservative.
